\scanpage{17}\mspage{15}
Donc la définition des $\pi_1(f,g)$ \inserted{(si c'est bien \unsure{ça|sa})} se prolonge de façon naturelle aux chemins infinis (\unsure{ess\textsuperscript{t}|ess.} finis) arbitraires.

$2^{\mathrm o})$ Supposons que $M$ admette un objet à la fois initial et final $0_M$ ($M$ ``\unsure{ponctuée|pointée}''). Je dis qu'alors \deleted{\unsure{toutes}} les catégories $\underline{\mathrm{Ch}}_{\infty}^{\Sigma}(a,b)$ sont \unsure{ponctuées|pointées} de façon naturelle : à savoir, par le chemin infini de $a$ à $b$ défini \inserted{(comme plus haut)} à partir de $f=0_{a,b}:a\unsure{\text{---}|\longrightarrow}b$, l'application nulle de $a$ à $b$, On \deleted{peut} a envie alors de regarder les $\underline{\mathrm{Ch}}_{\infty}^{\Sigma}(a,b)$ comme définissant des éléments dans le $\mathrm{Hot}_{\bullet}$ (``Hot \unsure{ponctué|pointé}''). C'est plus fin que de dire qu'on se donne une classe neutre
\[
e\longrightarrow\mathrm{Ch}(a,b)\qquad\text{dans Hot}
\]
car le foncteur de $\mathrm{Hot}_{\bullet}$ dans $\mathrm{Hot}$ (ou dans $e\backslash\mathrm{Hot}$) n'est pas \unsure{fidèle}\unsure{?|.} Le problème, c'est de définir les flèches
\[
\mathrm{Ch}(a,b)\times\mathrm{Ch}(b,c)\ \text{---}\ \mathrm{Ch}(a,c)
\]
dans $\mathrm{Hot}_{\bullet}$, et pas seulement dans $\mathrm{Hot}$, et avec associativité \inserted{(et les unités)} dans $\mathrm{Hot}_{\bullet}$ également. Si on applique la définition avec $\mathrm{Ch}^{-}(\unsure{a},b)$ et $\mathrm{Ch}^{+}(b,c)$, le composé des éléments nuls
\[
\begin{aligned}
&\unsure{a}\overset{0}{\text{---}}b\ \deleted{\unsure{\longleftarrow b\longrightarrow c}}\quad\text{de }\mathrm{Ch}^{\deleted{\unsure{-|+}}}(a,b)\quad\text{et}\\
&\deleted{\unsure{b\longrightarrow b\longrightarrow c}}\inserted{b\xleftarrow{\mathrm{id}}b\xrightarrow{0}c}\quad\text{de }\mathrm{Ch}^{+}(b,c)
\end{aligned}
\]
% archive sheet number: 16

